\documentclass[UTF8]{ctexart}

\usepackage[a4paper,margin=1in]{geometry}
\usepackage{amsmath,amssymb,mathtools,bm}
\usepackage{booktabs,longtable,array}
\usepackage{hyperref}
\usepackage{enumitem}

\hypersetup{
    colorlinks=true,
    linkcolor=black,
    urlcolor=blue,
    citecolor=black
}

\setlist[itemize]{leftmargin=2em}
\setlist[enumerate]{leftmargin=2em}

\title{《A Soft Sensor for Multirate Quality Variables Based on MC-CNN》方法报告}
\author{}
\date{2026-03-27}

\begin{document}

\maketitle
\tableofcontents
\newpage

\section{文档目标}

本文围绕 Song 等人在 2025 年发表的 MC-CNN（Multitemporal Channel CNN）方法，给出一份面向实现和理解的中文技术报告。报告重点回答以下问题：
\begin{itemize}
    \item 论文到底提出了什么方法，它想解决什么问题；
    \item 原始输入、训练标签、最终输出各自是什么，形状如何；
    \item 共享网络和并行预测网络的逐层数据流是什么样子；
    \item 论文中的“修改版 BP”在数学上等价于什么；
    \item 多时间通道为什么能提高稀疏质量变量的利用率。
\end{itemize}

文中会尽量保留论文原始记号，但会同时改写成更适合深度学习实现的张量视角。需要先说明一点：论文在网络结构部分给出了层维度，但没有像现代深度学习论文那样把每一层的张量轴定义得非常显式，因此本文对\textbf{卷积轴方向和若干中间张量形状}做了工程化推断。这些推断是有依据的，且能与论文给出的层维度完全对齐。

\section{方法总览}

\subsection{一句话概括}

MC-CNN 的核心思想是：\textbf{把最近 $M$ 个时刻的过程变量堆叠成 $M$ 个时间通道，用共享网络学习“过程变量到多时刻质量变量窗口”的映射，再用并行子网络把共享特征变成每个质量变量当前时刻的预测值}。同时，通过掩码式反向传播，让未采样时刻产生的空白标签不参与梯度计算。

\subsection{它解决的是什么问题}

工业软测量里经常出现如下情况：
\begin{itemize}
    \item 过程变量 $X$ 可以高频在线采样；
    \item 质量变量 $Y$ 只能低频采样，甚至不同质量变量采样周期还不同；
    \item 目标却是希望在每个时刻都输出质量变量的预测值。
\end{itemize}

如果直接用普通的 MIMO 回归模型，会遇到两个问题：
\begin{itemize}
    \item 大量时刻没有真实质量标签，监督信号过于稀疏；
    \item 只看“当前时刻”输入，难以学习过程变量和质量变量之间的时滞关系。
\end{itemize}

MC-CNN 针对这两个问题同时设计了两个机制：
\begin{itemize}
    \item \textbf{多时间通道}：把连续 $M$ 个时刻的过程变量打包输入，从而显式建模时序依赖；
    \item \textbf{修改版 BP}：对空白标签不回传梯度，使模型可以直接在带缺测标签的窗口上训练。
\end{itemize}

\subsection{与普通 MIMO-CNN 的区别}

论文里专门构造了一个对照模型 MIMO-CNN。两者最本质的区别不是“有没有 CNN”，而是\textbf{最后一层监督方式不同}：
\begin{itemize}
    \item MIMO-CNN 只用单时刻标签训练，本质上仍然主要学习“当前输入到当前输出”；
    \item MC-CNN 用 $M$ 个时刻组成的质量变量窗口 $D_y$ 监督共享网络，因此共享特征会被迫对多时刻标签同时有解释力。
\end{itemize}

这一点是整篇论文最值得抓住的设计核心。

\section{符号、数据形式与张量形状}

\subsection{基础符号}

设：
\begin{itemize}
    \item $T$ 为总采样时刻数；
    \item $m$ 为过程变量个数；
    \item $n$ 为质量变量个数；
    \item $M$ 为多时间通道数，也就是时间窗口长度；
    \item $x_t \in \mathbb{R}^{m}$ 为第 $t$ 个时刻的过程变量向量；
    \item $y_t \in \mathbb{R}^{n}$ 为第 $t$ 个时刻的质量变量向量。
\end{itemize}

由于质量变量是多采样率的，所以 $y_t$ 的某些分量在很多时刻并不存在真实观测值。论文把这些位置视为 blank value（空白值）。

\subsection{标准化和空白值处理后的原始矩阵}

论文在预处理后得到
\begin{equation}
    X^{nb} = [x_1, x_2, \dots, x_T]^\top \in \mathbb{R}^{T \times m},
\end{equation}
\begin{equation}
    Y^{nb} = [y_1, y_2, \dots, y_T]^\top \in \mathbb{R}^{T \times n}.
\end{equation}

这里上标 $nb$ 可以理解为“完成归一化并处理 blank 值之后”的数据形式。严格来说：
\begin{itemize}
    \item $X^{nb}$ 一般是完整观测的；
    \item $Y^{nb}$ 中未采样位置会被填上一个特殊值，训练时再通过修改版 BP 忽略掉。
\end{itemize}

从现代实现角度看，更推荐把这个“特殊值”显式写成掩码矩阵
\begin{equation}
    R \in \{0,1\}^{T \times n},
\end{equation}
其中 $R_{t,q}=1$ 表示第 $t$ 个时刻第 $q$ 个质量变量有真实观测。

\subsection{多时间通道窗口化}

论文把长度为 $M$ 的滑动窗口拼接成三维张量。若记
\begin{equation}
    S = T - M + 1,
\end{equation}
则有
\begin{equation}
    D_x =
    \begin{bmatrix}
        [x_1, x_2, \dots, x_M]^\top \\
        [x_2, x_3, \dots, x_{M+1}]^\top \\
        \vdots \\
        [x_{T-M+1}, x_{T-M+2}, \dots, x_T]^\top
    \end{bmatrix}
    \in \mathbb{R}^{S \times M \times m},
\end{equation}
\begin{equation}
    D_y =
    \begin{bmatrix}
        [y_1, y_2, \dots, y_M]^\top \\
        [y_2, y_3, \dots, y_{M+1}]^\top \\
        \vdots \\
        [y_{T-M+1}, y_{T-M+2}, \dots, y_T]^\top
    \end{bmatrix}
    \in \mathbb{R}^{S \times M \times n}.
\end{equation}

因此，\textbf{单个训练样本}的形状是：
\begin{equation}
    d_x^{(s)} \in \mathbb{R}^{M \times m}, \qquad
    d_y^{(s)} \in \mathbb{R}^{M \times n}.
\end{equation}

若按 batch 训练，batch size 为 $B$，则输入输出张量分别是：
\begin{equation}
    \mathcal{X}_{\text{batch}} \in \mathbb{R}^{B \times M \times m}, \qquad
    \mathcal{Y}_{\text{batch}} \in \mathbb{R}^{B \times M \times n}.
\end{equation}

\subsection{一个非常重要的直觉}

窗口 $D_x$ 与 $D_y$ 的构造意味着：\textbf{模型训练时不是在看一个时刻，而是在看一段长度为 $M$ 的局部时间片段}。这让共享网络有机会学习：
\begin{itemize}
    \item 哪个质量变量更依赖最近时刻的过程变量；
    \item 哪个质量变量更依赖更早的滞后信息；
    \item 多个质量变量之间在同一时间片上的互补模式。
\end{itemize}

\section{MC-CNN 的整体结构}

\subsection{结构分成两部分}

论文把 MC-CNN 明确分成两块：
\begin{itemize}
    \item \textbf{共享网络（shared network）}：输入是 $D_x$，监督标签是 $D_y$，用于提取共享特征 $h(k)$；
    \item \textbf{并行预测网络（parallel prediction network）}：每个质量变量各有一个子网络，以共享特征 $h(k)$ 为输入，输出该质量变量当前时刻的预测值。
\end{itemize}

图示上可以把它理解为：
\begin{equation}
    D_x \longrightarrow \text{Shared Network} \longrightarrow h(k)
    \longrightarrow
    \begin{cases}
        \text{Head}_1 \to \hat{y}_1(k) \\
        \text{Head}_2 \to \hat{y}_2(k) \\
        \vdots \\
        \text{Head}_n \to \hat{y}_n(k)
    \end{cases}
\end{equation}

同时，共享网络内部还会生成一个窗口级输出 $\hat{D}_y \in \mathbb{R}^{M \times n}$，用于第一阶段的预训练。

\subsection{共享网络的角色}

共享网络的目标不是直接给出最终在线预测，而是先回答一个更容易学的问题：
\begin{quote}
给定最近 $M$ 个时刻的过程变量窗口，能否恢复同一窗口里的质量变量模式？
\end{quote}

这样做有两个好处：
\begin{itemize}
    \item 共享层能被多时刻监督，得到更稳定的时序特征；
    \item 稀疏标签会在多个窗口位置重复出现，提高有效监督次数。
\end{itemize}

\subsection{并行预测网络的角色}

并行预测网络才是最终软测量输出层。它把共享特征 $h(k)$ 投影为当前时刻各质量变量的预测值：
\begin{equation}
    \hat{y}(k) = [\hat{y}_1(k), \hat{y}_2(k), \dots, \hat{y}_n(k)]^\top \in \mathbb{R}^{n}.
\end{equation}

因此，从任务分工上看：
\begin{itemize}
    \item 共享网络负责“抽取多时刻、多变量共同有用的表示”；
    \item 每个预测头负责“把共享表示变成某个具体质量变量的标量回归结果”。
\end{itemize}

\section{逐层张量数据流分析}

\subsection{先说一个关键推断}

论文给出的 TE 工况共享网络层维度为：
\begin{equation}
    [[5,33]-[5,40]-[3,38]-[36]-[25]-[5,2]].
\end{equation}

这组维度几乎唯一地说明了如下张量解释：
\begin{itemize}
    \item 第一维的 $5$ 不是 batch，也不是特征数，而是\textbf{多时间通道数} $M$；
    \item 第一层全连接是对每个时间通道上的 $33$ 维过程变量向量独立做线性变换，得到 $40$ 维；
    \item 之后的卷积层把 $M=5$ 当作输入通道数，把投影后的特征轴长度 $40$ 当作卷积的一维长度。
\end{itemize}

也就是说，\textbf{这不是“沿时间轴滑动的时域卷积”，而是“时间滞后作为通道，卷积沿投影特征轴滑动”}。时间信息是通过通道混合进入卷积的，这一点非常重要。

\subsection{共享网络：一般形状}

设一个 batch 的输入为
\begin{equation}
    X_b \in \mathbb{R}^{B \times M \times m}.
\end{equation}

共享网络可以写成以下数据流：

\paragraph{步骤 1：逐通道全连接投影}
\begin{equation}
    H^{(0)} = \phi(X_b W_1 + b_1),
\end{equation}
其中 $\phi$ 表示 SELU 激活。若投影后宽度为 $d_1$，则
\begin{equation}
    H^{(0)} \in \mathbb{R}^{B \times M \times d_1}.
\end{equation}

\paragraph{步骤 2：第一层 1D 卷积}
\begin{equation}
    H^{(1)} = \phi(\operatorname{Conv1d}(H^{(0)})),
\end{equation}
若输出通道数为 $c_1$、卷积核宽度为 $k_1$、步长为 1，则
\begin{equation}
    H^{(1)} \in \mathbb{R}^{B \times c_1 \times (d_1-k_1+1)}.
\end{equation}

\paragraph{步骤 3：第二层 1D 卷积}
\begin{equation}
    H^{(2)} = \phi(\operatorname{Conv1d}(H^{(1)})),
\end{equation}
若第二层输出通道数为 $1$、卷积核宽度为 $k_2$，则
\begin{equation}
    H^{(2)} \in \mathbb{R}^{B \times 1 \times (d_1-k_1-k_2+2)}.
\end{equation}

\paragraph{步骤 4：展平}
\begin{equation}
    z = \operatorname{Flatten}(H^{(2)}) \in \mathbb{R}^{B \times (d_1-k_1-k_2+2)}.
\end{equation}

\paragraph{步骤 5：共享隐特征}
\begin{equation}
    h = \phi(z W_h + b_h) \in \mathbb{R}^{B \times d_h}.
\end{equation}

\paragraph{步骤 6：窗口级输出}
\begin{equation}
    \hat{D}_y = \operatorname{reshape}(h W_o + b_o) \in \mathbb{R}^{B \times M \times n}.
\end{equation}

\subsection{TE 工况的具体形状}

在 TE 实验中，论文给出：
\begin{itemize}
    \item $M=5$；
    \item $m=33$；
    \item $n=2$；
    \item 共享网络维度为 $[[5,33]-[5,40]-[3,38]-[36]-[25]-[5,2]]$。
\end{itemize}

于是 batch 级形状可以明确写成：
\begin{align}
    X_b &\in \mathbb{R}^{B \times 5 \times 33}, \\
    H^{(0)} &\in \mathbb{R}^{B \times 5 \times 40}, \\
    H^{(1)} &\in \mathbb{R}^{B \times 3 \times 38}, \\
    H^{(2)} &\in \mathbb{R}^{B \times 1 \times 36}, \\
    z &\in \mathbb{R}^{B \times 36}, \\
    h &\in \mathbb{R}^{B \times 25}, \\
    \hat{D}_y &\in \mathbb{R}^{B \times 5 \times 2}.
\end{align}

这里的两个卷积核宽度都可由维度反推出为 3，因为：
\begin{equation}
    40 \xrightarrow{k=3} 38 \xrightarrow{k=3} 36.
\end{equation}

\subsection{储能数据的具体形状}

在储能数据实验中，论文给出的共享网络维度为
\begin{equation}
    [[5,640]-[5,320]-[5,160]-[5,80]-[5,40]-[3,38]-[36]-[25]-[5,2]].
\end{equation}

因此对应的数据流是：
\begin{align}
    X_b &\in \mathbb{R}^{B \times 5 \times 640}, \\
    H^{(0)}_1 &\in \mathbb{R}^{B \times 5 \times 320}, \\
    H^{(0)}_2 &\in \mathbb{R}^{B \times 5 \times 160}, \\
    H^{(0)}_3 &\in \mathbb{R}^{B \times 5 \times 80}, \\
    H^{(0)}_4 &\in \mathbb{R}^{B \times 5 \times 40}, \\
    H^{(1)} &\in \mathbb{R}^{B \times 3 \times 38}, \\
    H^{(2)} &\in \mathbb{R}^{B \times 1 \times 36}, \\
    z &\in \mathbb{R}^{B \times 36}, \\
    h &\in \mathbb{R}^{B \times 25}, \\
    \hat{D}_y &\in \mathbb{R}^{B \times 5 \times 2}.
\end{align}

可以看出，两组实验虽然输入维度差别很大，但作者都先把每个时间通道上的高维过程变量压到 40 维，再进入相同的卷积--展平--隐层结构。

\subsection{并行预测网络的形状}

论文给出的每个质量变量预测头结构是：
\begin{equation}
    [25-5-1].
\end{equation}

这意味着第 $q$ 个质量变量的预测头可以写成
\begin{equation}
    \hat{y}_q = W_{q,2}\,\phi(W_{q,1} h + b_{q,1}) + b_{q,2},
\end{equation}
其中
\begin{equation}
    h \in \mathbb{R}^{B \times 25},
    \qquad
    \hat{y}_q \in \mathbb{R}^{B \times 1}.
\end{equation}

把所有预测头拼起来之后：
\begin{equation}
    \hat{Y}_{\text{latest}} \in \mathbb{R}^{B \times n}.
\end{equation}

如果是论文中的两个质量变量案例，则
\begin{equation}
    \hat{Y}_{\text{latest}} \in \mathbb{R}^{B \times 2}.
\end{equation}

\subsection{共享输出和最终输出不是一回事}

这是阅读论文时最容易混淆的一点。

\paragraph{共享网络输出}
\begin{equation}
    \hat{D}_y \in \mathbb{R}^{B \times M \times n}
\end{equation}
是窗口级输出，用来做第一阶段训练监督。

\paragraph{并行预测网络输出}
\begin{equation}
    \hat{Y}_{\text{latest}} \in \mathbb{R}^{B \times n}
\end{equation}
才是最终用于在线软测量的当前时刻预测。

所以更准确地说，MC-CNN 有\textbf{两个输出层次}：
\begin{itemize}
    \item 一个是“恢复整个质量窗口”的辅助输出；
    \item 一个是“预测当前时刻各质量变量”的最终输出。
\end{itemize}

\section{卷积层到底在做什么}

\subsection{卷积不是直接沿时间轴滑动}

如果采用 PyTorch/深度学习框架的标准约定，1D 卷积的输入通常是
\begin{equation}
    \mathbb{R}^{B \times C \times L},
\end{equation}
其中 $C$ 是通道数，$L$ 是卷积长度。

在本文结构里，最自然的解释是：
\begin{itemize}
    \item $C=M$，即 $M$ 个时间通道；
    \item $L=d_1$，即每个时间通道经全连接投影后的特征长度。
\end{itemize}

因此第一层卷积的一个输出元素可以写成
\begin{equation}
    H^{(1)}_{b,r,p}
    =
    \phi\!\left(
        \sum_{c=1}^{M}
        \sum_{\delta=0}^{k_1-1}
        K^{(1)}_{r,c,\delta}\,
        H^{(0)}_{b,c,p+\delta}
        + b^{(1)}_r
    \right),
\end{equation}
其中：
\begin{itemize}
    \item $r$ 是输出卷积核编号；
    \item $c$ 遍历所有时间通道；
    \item $\delta$ 遍历卷积核局部感受野。
\end{itemize}

这个式子说明：\textbf{每个卷积核会同时混合所有时间通道的信息}。因此，虽然卷积核不是沿“时间轴”平移，但时间滞后信息已经通过通道维进入了卷积运算。

\subsection{为什么这样仍然能学习时滞关系}

原因在于：第 $1,\dots,M$ 个通道本身就分别代表不同的相对时间位置。比如当 $M=5$ 时：
\begin{equation}
    d_x^{(s)} =
    \begin{bmatrix}
        x_s \\
        x_{s+1} \\
        x_{s+2} \\
        x_{s+3} \\
        x_{s+4}
    \end{bmatrix}.
\end{equation}

模型看到的并不是一组无序向量，而是一个带有固定相对顺序的通道集合。不同通道权重不同，就意味着模型能区分：
\begin{itemize}
    \item 当前质量值更依赖 $x_t$；
    \item 还是更依赖 $x_{t-1}$、$x_{t-2}$；
    \item 甚至不同质量变量依赖的滞后步长不同。
\end{itemize}

\subsection{没有池化层的含义}

论文明确说 MC-CNN 不使用 pooling 层，因为作者担心池化会滤掉“次要特征”，而这些次要特征对软测量回归可能仍然有用。这个选择说明作者更关心\textbf{保留细粒度回归信息}，而不是像图像分类那样追求平移不变性。

\section{修改版 BP 的数学本质}

\subsection{论文写法}

论文定义了一个函数
\begin{equation}
    g(x-y)=
    \begin{cases}
        x-y, & y \in y_r, \\
        0, & y \notin y_r,
    \end{cases}
\end{equation}
其中 $y_r$ 表示真实采样值集合。然后把共享网络的损失写成
\begin{equation}
    J_s = \frac{1}{2 N Z} \sum_{i=1}^{N}\sum_{j=1}^{M}
    \left\{
        g\!\left(\tilde{y}(i,j)-y(i,j)\right)
    \right\}^2,
\end{equation}
并把预测头损失写成
\begin{equation}
    J_p = \frac{1}{2 Z} \sum_{i=1}^{N}
    \left\{
        g\!\left(\tilde{y}(i)-y(i)\right)
    \right\}^2.
\end{equation}

这里 $Z$ 是标签里真实值的个数。

\subsection{更清晰的张量化改写}

更现代、也更容易实现的写法是直接引入 mask。设
\begin{equation}
    Q \in \{0,1\}^{B \times M \times n}
\end{equation}
表示窗口标签是否真实可用，那么共享网络损失可以写成
\begin{equation}
    \mathcal{L}_{\text{shared}}
    =
    \frac{1}{2 \sum Q}
    \left\|
        Q \odot (\hat{D}_y - D_y)
    \right\|_F^2.
\end{equation}

若定义当前时刻标签掩码
\begin{equation}
    q^{\text{latest}} \in \{0,1\}^{B \times n},
\end{equation}
则并行预测头的损失为
\begin{equation}
    \mathcal{L}_{\text{head}}
    =
    \frac{1}{2 \sum q^{\text{latest}}}
    \left\|
        q^{\text{latest}} \odot (\hat{Y}_{\text{latest}} - Y_{\text{latest}})
    \right\|_F^2.
\end{equation}

这里 $\odot$ 表示按元素乘法。

\subsection{为什么空白值不会参与反向传播}

对共享损失中的任意一个位置 $(b,c,q)$，有
\begin{equation}
    \frac{\partial \mathcal{L}_{\text{shared}}}{\partial \hat{D}_{y,b,c,q}}
    =
    \frac{Q_{b,c,q}}{\sum Q}
    \left(
        \hat{D}_{y,b,c,q} - D_{y,b,c,q}
    \right).
\end{equation}

如果该位置没有真实观测，那么 $Q_{b,c,q}=0$，于是
\begin{equation}
    \frac{\partial \mathcal{L}_{\text{shared}}}{\partial \hat{D}_{y,b,c,q}} = 0.
\end{equation}

这就说明：\textbf{未采样位置对损失和梯度都没有贡献}。论文所谓的“修改版 BP”，本质上就是今天常说的 \textbf{masked regression loss}。

\subsection{它和 dropout 的区别}

论文特别强调这和 dropout 不是一回事。这个判断是正确的：
\begin{itemize}
    \item dropout 是一种随机正则化，目的是防止过拟合；
    \item MC-CNN 这里屏蔽的是\textbf{数据中本来就不存在真实标签的位置}，目的是避免错误监督。
\end{itemize}

前者是人为随机丢弃，后者是对缺失监督做严格掩码，二者功能和数学语义都不同。

\section{为什么多时间通道能“重复利用”稀疏标签}

\subsection{论文里的直觉}

论文举例说，当 $M=5$ 时，一个真实标签最多可以放到输出层的 5 个不同位置，因此“一个真实值可以被利用五次”。这句话的意思不是复制标签，而是指\textbf{同一个真实采样时刻会出现在多个滑动窗口中}。

\subsection{精确地写出来}

设第 $t$ 个时刻第 $q$ 个质量变量有真实采样值 $y_{t,q}$。它会出现在所有满足
\begin{equation}
    s \le t \le s+M-1
\end{equation}
的窗口中，也就是
\begin{equation}
    s \in [t-M+1,\; t]
\end{equation}
与合法窗口起点区间 $[1,S]$ 的交集内。

因此，这个真实值参与训练的次数为
\begin{equation}
    c_t
    =
    \left|
        [t-M+1,\; t] \cap [1,S]
    \right|.
\end{equation}

对中间时刻的样本，通常有
\begin{equation}
    c_t = M.
\end{equation}

边界附近的样本参与次数会少一些，但\textbf{绝大多数内部样本都能被重复利用 $M$ 次}。

\subsection{这为什么有用}

因为每次它出现时，对应的输入上下文都不同。比如同一个真实值 $y_t$，可能分别与如下过程窗口配对：
\begin{equation}
    [x_{t-4},x_{t-3},x_{t-2},x_{t-1},x_t],
\end{equation}
\begin{equation}
    [x_{t-3},x_{t-2},x_{t-1},x_t,x_{t+1}],
\end{equation}
\begin{equation}
    \dots
\end{equation}

于是模型不只学到“当前过程状态和当前质量值的关系”，还学到“不同相对滞后位置下，哪些过程模式更能解释这个质量值”。这正是 MC-CNN 想捕捉的 temporal correlations。

\section{训练流程与推理流程}

\subsection{论文描述的训练流程}

根据论文正文，训练大致是一个两阶段过程：
\begin{enumerate}
    \item 收集过程变量和质量变量，划分训练集与测试集；
    \item 对数据做 Z-score 标准化；
    \item 对未采样质量变量填入特殊值，并在训练中通过修改版 BP 忽略；
    \item 把 $X$ 和 $Y$ 窗口化为 $D_x$ 与 $D_y$；
    \item 先用 $(D_x, D_y)$ 预训练共享网络，提取共享特征 $h(k)$；
    \item 再用共享特征 $h(k)$ 和当前时刻质量标签训练并行预测头；
    \item 推理时，把最近 $M$ 个时刻的过程变量输入模型，输出当前时刻各质量变量预测值。
\end{enumerate}

\subsection{共享网络预训练阶段的数据流}

在 batch 级别，这一阶段的输入输出可以总结为：
\begin{equation}
    \mathbb{R}^{B \times M \times m}
    \longrightarrow
    \mathbb{R}^{B \times M \times n}.
\end{equation}

也就是
\begin{equation}
    X_b \xrightarrow{\text{shared network}} \hat{D}_y
    \xrightarrow{\text{masked loss}} D_y.
\end{equation}

这一步训练出来的重点不是最终输出层，而是中间特征
\begin{equation}
    h \in \mathbb{R}^{B \times 25}.
\end{equation}

\subsection{并行预测头训练阶段的数据流}

第二阶段对每个质量变量单独训练一个头。若批量当前时刻标签记为
\begin{equation}
    Y_{\text{latest}} \in \mathbb{R}^{B \times n},
\end{equation}
则这一阶段是
\begin{equation}
    h \in \mathbb{R}^{B \times 25}
    \longrightarrow
    \hat{Y}_{\text{latest}} \in \mathbb{R}^{B \times n}
    \xrightarrow{\text{masked loss}}
    Y_{\text{latest}}.
\end{equation}

\paragraph{一个实现细节}
论文文字表述更像是“先训练共享网络，再训练预测头”。它没有明确说明第二阶段是否继续联合更新共享网络参数。最保守的理解是：
\begin{itemize}
    \item 第一阶段学习共享特征；
    \item 第二阶段把共享特征当作固定输入训练各预测头。
\end{itemize}

如果采用端到端联合微调，那属于与论文主线一致、但略偏工程扩展的实现方式。

\subsection{在线推理的数据流}

在线推理时，不需要真实质量标签。只要最新的过程变量窗口可得，即可执行：
\begin{equation}
    [x_{t-M+1}, \dots, x_t]
    \longrightarrow
    h(t)
    \longrightarrow
    [\hat{y}_1(t), \dots, \hat{y}_n(t)].
\end{equation}

因此最终软测量器的输入输出关系是：
\begin{equation}
    \mathbb{R}^{M \times m} \to \mathbb{R}^{n}.
\end{equation}

这也解释了为什么论文说 MC-CNN 可以在\textbf{高频}上输出质量变量预测值：即使真实质量变量只有低频采样，模型也可以利用每个时刻都有的过程变量持续推断。

\section{实例化到论文中的两个实验}

\subsection{TE 过程}

论文中 TE 过程的设置为：
\begin{itemize}
    \item 过程变量数 $m=33$；
    \item 质量变量数 $n=2$；
    \item 多时间通道数 $M=5$；
    \item 共享网络维度：$[[5,33]-[5,40]-[3,38]-[36]-[25]-[5,2]]$；
    \item 每个质量变量预测头：$[25-5-1]$；
    \item 超参数表给出：预训练学习率 $0.01$、epoch $8$、batch size $32$；微调学习率 $0.01$、epoch $6$、batch size $32$。
\end{itemize}

因此 TE 过程里，一个 batch 的典型前向形状是
\begin{equation}
    B \times 5 \times 33
    \to
    B \times 5 \times 40
    \to
    B \times 3 \times 38
    \to
    B \times 1 \times 36
    \to
    B \times 36
    \to
    B \times 25
    \to
    B \times 5 \times 2,
\end{equation}
随后再由并行头得到
\begin{equation}
    B \times 25 \to B \times 2.
\end{equation}

\subsection{储能数据}

论文中储能数据的设置为：
\begin{itemize}
    \item 过程变量数 $m=640$；
    \item 质量变量数 $n=2$；
    \item 多时间通道数 $M=5$；
    \item 共享网络维度：$[[5,640]-[5,320]-[5,160]-[5,80]-[5,40]-[3,38]-[36]-[25]-[5,2]]$；
    \item 每个质量变量预测头：$[25-5-1]$；
    \item 超参数表给出：预训练学习率 $0.01$、epoch $6$、batch size $32$；微调学习率 $0.01$、epoch $6$、batch size $32$。
\end{itemize}

这里的主要区别只是前端增加了若干层逐通道全连接，用来把 $640$ 维高维输入逐步压缩到 $40$ 维。后面的卷积、展平和隐层结构与 TE 设置保持一致。

\section{从机器学习原理看 MC-CNN}

\subsection{它本质上是在做什么}

如果不拘泥于论文术语，MC-CNN 可以理解为下面这个组合：
\begin{itemize}
    \item \textbf{滑动窗口监督学习}：把时序问题转成局部窗口回归问题；
    \item \textbf{多任务学习}：共享网络同时服务于多个质量变量；
    \item \textbf{带掩码的缺失标签回归}：未观测标签不计入损失；
    \item \textbf{滞后建模}：最近 $M$ 个时刻被显式编码到输入通道中。
\end{itemize}

\subsection{为什么共享网络有效}

共享网络被要求同时解释：
\begin{itemize}
    \item 同一个质量变量在不同时间位置上的值；
    \item 不同质量变量之间的联合模式。
\end{itemize}

因此学到的隐特征 $h(k)$ 不太可能只对单个输出有用，它更像是一个\textbf{对局部过程状态的多变量时序摘要}。这也是后续每个质量变量预测头都能复用它的原因。

\subsection{为什么适合多采样率质量变量}

对多采样率任务，最麻烦的是监督很稀疏。MC-CNN 的做法不是先插值补齐所有质量变量，而是：
\begin{itemize}
    \item 保留真实采样时刻；
    \item 用 mask 屏蔽未采样位置；
    \item 通过滑动窗口让真实采样值出现在多个训练样本中。
\end{itemize}

这比直接插值更稳妥，因为它没有把“伪标签”当成真值强行喂给网络。

\section{方法的优点、边界与实现注意点}

\subsection{优点}

\begin{itemize}
    \item 能同时处理多个质量变量；
    \item 能处理不同质量变量采样周期不同的情况；
    \item 通过窗口监督提高低频标签利用率；
    \item 结构较小，作者强调参数量和训练时间都比较可控；
    \item 推理阶段只依赖过程变量窗口，适合在线预测。
\end{itemize}

\subsection{阅读论文时容易误解的地方}

\begin{itemize}
    \item “多时间通道 CNN”不等价于标准 TCN。它不是在时间轴上做因果卷积，而是把时间滞后编码到通道中。
    \item $[5,2]$ 这个输出维度是共享网络窗口输出，不是最终在线输出。最终在线输出是 $[2]$。
    \item 论文的“修改版 BP”本质上就是 masked MSE，只是用特殊值加条件函数的方式表达。
\end{itemize}

\subsection{方法边界}

\begin{itemize}
    \item 采样间隔若非常不规则，固定长度窗口 $M$ 的表达能力会受限；
    \item 质量变量极少、过程变量极多时，前端全连接层仍然可能有较大参数量；
    \item 该结构更像“滞后特征提取器”，并不像 Transformer 或 RNN 那样显式建模长程依赖；
    \item 论文把 $M$ 当作重要超参数，选太小学不到足够时滞，选太大则训练开销上升。
\end{itemize}

\section{结论}

MC-CNN 的真正创新点不在于“把 CNN 用到软测量”本身，而在于它把\textbf{多时间通道、窗口级监督、并行多输出预测、以及掩码式反向传播}组合到了一起，专门应对“过程变量高频、质量变量低频且多采样率”的场景。

如果用一句更工程化的话来总结，它做的是：
\begin{quote}
先用最近 $M$ 个时刻的过程变量去重建一个带缺测的质量变量窗口，从中学到共享时序特征；再把这个共享特征送入多个小回归头，输出每个质量变量当前时刻的高频预测值。
\end{quote}

从张量角度看，MC-CNN 的关键数据流更准确地应写成
\begin{equation}
    \mathbb{R}^{B \times M \times m}
    \to
    \mathbb{R}^{B \times 25}
    \to
    \begin{cases}
        \mathbb{R}^{B \times M \times n}, \\
        \mathbb{R}^{B \times n},
    \end{cases}
\end{equation}
其中共享隐特征先由窗口输入得到，再分别通向窗口级监督输出和最终在线软测量输出。

从数学角度看，它的“修改版 BP”可以非常清楚地理解为：
\begin{quote}
对真实采样标签做均方误差，对空白标签位置乘零掩码，不让错误监督传播到网络参数。
\end{quote}

这也是 MC-CNN 能在低频质量变量条件下仍然保持较好训练效果的关键所在。

\end{document}
